% ══════════════════════════════════════════════════════════════════════════════
%  المشترك — التمهيد الموحّد وطبقة التصميم لوثائق الدروس (article, عربي RTL)
%  يُصرَّف بـ XeLaTeX حصراً.
%
%  يُستدعى بـ (من مجلد الدرس):
%    \input{"../../shared/common.tex"}
%
%  يوفّر: التمهيد، الخطوط، الألوان (أزرق/أسود/أبيض)، إطار الصفحة، الترويسة،
%          التذييل (الصفحة س/ص)، شبكة الجداول بخطوط عمودية، صناديق وشارات وأيقونات،
%          وكل ماكروهات الوثائق عبر design.tex + lesson-macros.tex.
% ══════════════════════════════════════════════════════════════════════════════

% ── Geometry ──────────────────────────────────────────────────────────────────
\usepackage[
  top=2.8cm,
  bottom=2.4cm,
  left=2.2cm,
  right=2.2cm,
  headheight=26pt          % MUST be >= 26pt (fancyhdr)
]{geometry}

% ── Packages: fancyhdr and friends BEFORE polyglossia ────────────────────────
\usepackage{fancyhdr}
\usepackage{titlesec}
\usepackage{enumitem}
\usepackage[table]{xcolor}   % table → \rowcolors / \rowcolor / \cellcolor
\usepackage{tikz}
\usetikzlibrary{positioning}
\usetikzlibrary{arrows.meta}
\usetikzlibrary{shapes.geometric}
\usepackage{eso-pic}
\usepackage{tcolorbox}
\tcbuselibrary{skins, breakable}
\usepackage{array}
\usepackage{tabularx}
\usepackage{longtable}
\usepackage{multirow}
\usepackage{lastpage}      % للتذييل: الصفحة س / ص
\usepackage{amsmath}
\usepackage{etoolbox}
\usepackage{hyperref}

% روابط بلا إطار (يمنع المستطيل حول رقم الصفحة في التذييل)
\hypersetup{hidelinks}

% ── Arabic / RTL — load AFTER fancyhdr ───────────────────────────────────────
\usepackage{polyglossia}
\setmainlanguage[numerals=western]{arabic}
\setotherlanguage{english}

% ── Fonts (Amiri — خط النسخ المعتمد) ────────────────────────────────────────
\setmainfont{Amiri}[
  BoldFont = * Bold,
  ItalicFont = * Italic,
]
\newfontfamily\arabicfont{Amiri}[
  Script = Arabic,
  BoldFont = * Bold,
]
\setmonofont{DejaVu Sans Mono}

% ── TYPOGRAPHY: تدرّج طباعي وتهوية ─────────────────────────────────────────
\linespread{1.18}
\setlength{\parindent}{0pt}
\setlength{\parskip}{5pt plus 1pt}
\renewcommand{\arraystretch}{1.2}


% ── مسار مجلد shared (قابل للتجاوز قبل % ══════════════════════════════════════════════════════════════════════════════
%  المشترك — التمهيد الموحّد وطبقة التصميم لوثائق الدروس (article, عربي RTL)
%  يُصرَّف بـ XeLaTeX حصراً.
%
%  يُستدعى بـ (من مجلد الدرس):
%    \input{"../../shared/common.tex"}
%
%  يوفّر: التمهيد، الخطوط، الألوان (أزرق/أسود/أبيض)، إطار الصفحة، الترويسة،
%          التذييل (الصفحة س/ص)، شبكة الجداول بخطوط عمودية، صناديق وشارات وأيقونات،
%          وكل ماكروهات الوثائق عبر design.tex + lesson-macros.tex.
% ══════════════════════════════════════════════════════════════════════════════

% ── Geometry ──────────────────────────────────────────────────────────────────
\usepackage[
  top=2.8cm,
  bottom=2.4cm,
  left=2.2cm,
  right=2.2cm,
  headheight=26pt          % MUST be >= 26pt (fancyhdr)
]{geometry}

% ── Packages: fancyhdr and friends BEFORE polyglossia ────────────────────────
\usepackage{fancyhdr}
\usepackage{titlesec}
\usepackage{enumitem}
\usepackage[table]{xcolor}   % table → \rowcolors / \rowcolor / \cellcolor
\usepackage{tikz}
\usetikzlibrary{positioning}
\usetikzlibrary{arrows.meta}
\usetikzlibrary{shapes.geometric}
\usepackage{eso-pic}
\usepackage{tcolorbox}
\tcbuselibrary{skins, breakable}
\usepackage{array}
\usepackage{tabularx}
\usepackage{longtable}
\usepackage{multirow}
\usepackage{lastpage}      % للتذييل: الصفحة س / ص
\usepackage{amsmath}
\usepackage{etoolbox}
\usepackage{hyperref}

% روابط بلا إطار (يمنع المستطيل حول رقم الصفحة في التذييل)
\hypersetup{hidelinks}

% ── Arabic / RTL — load AFTER fancyhdr ───────────────────────────────────────
\usepackage{polyglossia}
\setmainlanguage[numerals=western]{arabic}
\setotherlanguage{english}

% ── Fonts (Amiri — خط النسخ المعتمد) ────────────────────────────────────────
\setmainfont{Amiri}[
  BoldFont = * Bold,
  ItalicFont = * Italic,
]
\newfontfamily\arabicfont{Amiri}[
  Script = Arabic,
  BoldFont = * Bold,
]
\setmonofont{DejaVu Sans Mono}

% ── TYPOGRAPHY: تدرّج طباعي وتهوية ─────────────────────────────────────────
\linespread{1.18}
\setlength{\parindent}{0pt}
\setlength{\parskip}{5pt plus 1pt}
\renewcommand{\arraystretch}{1.2}


% ── مسار مجلد shared (قابل للتجاوز قبل % ══════════════════════════════════════════════════════════════════════════════
%  المشترك — التمهيد الموحّد وطبقة التصميم لوثائق الدروس (article, عربي RTL)
%  يُصرَّف بـ XeLaTeX حصراً.
%
%  يُستدعى بـ (من مجلد الدرس):
%    \input{"../../shared/common.tex"}
%
%  يوفّر: التمهيد، الخطوط، الألوان (أزرق/أسود/أبيض)، إطار الصفحة، الترويسة،
%          التذييل (الصفحة س/ص)، شبكة الجداول بخطوط عمودية، صناديق وشارات وأيقونات،
%          وكل ماكروهات الوثائق عبر design.tex + lesson-macros.tex.
% ══════════════════════════════════════════════════════════════════════════════

% ── Geometry ──────────────────────────────────────────────────────────────────
\usepackage[
  top=2.8cm,
  bottom=2.4cm,
  left=2.2cm,
  right=2.2cm,
  headheight=26pt          % MUST be >= 26pt (fancyhdr)
]{geometry}

% ── Packages: fancyhdr and friends BEFORE polyglossia ────────────────────────
\usepackage{fancyhdr}
\usepackage{titlesec}
\usepackage{enumitem}
\usepackage[table]{xcolor}   % table → \rowcolors / \rowcolor / \cellcolor
\usepackage{tikz}
\usetikzlibrary{positioning}
\usetikzlibrary{arrows.meta}
\usetikzlibrary{shapes.geometric}
\usepackage{eso-pic}
\usepackage{tcolorbox}
\tcbuselibrary{skins, breakable}
\usepackage{array}
\usepackage{tabularx}
\usepackage{longtable}
\usepackage{multirow}
\usepackage{lastpage}      % للتذييل: الصفحة س / ص
\usepackage{amsmath}
\usepackage{etoolbox}
\usepackage{hyperref}

% روابط بلا إطار (يمنع المستطيل حول رقم الصفحة في التذييل)
\hypersetup{hidelinks}

% ── Arabic / RTL — load AFTER fancyhdr ───────────────────────────────────────
\usepackage{polyglossia}
\setmainlanguage[numerals=western]{arabic}
\setotherlanguage{english}

% ── Fonts (Amiri — خط النسخ المعتمد) ────────────────────────────────────────
\setmainfont{Amiri}[
  BoldFont = * Bold,
  ItalicFont = * Italic,
]
\newfontfamily\arabicfont{Amiri}[
  Script = Arabic,
  BoldFont = * Bold,
]
\setmonofont{DejaVu Sans Mono}

% ── TYPOGRAPHY: تدرّج طباعي وتهوية ─────────────────────────────────────────
\linespread{1.18}
\setlength{\parindent}{0pt}
\setlength{\parskip}{5pt plus 1pt}
\renewcommand{\arraystretch}{1.2}


% ── مسار مجلد shared (قابل للتجاوز قبل % ══════════════════════════════════════════════════════════════════════════════
%  المشترك — التمهيد الموحّد وطبقة التصميم لوثائق الدروس (article, عربي RTL)
%  يُصرَّف بـ XeLaTeX حصراً.
%
%  يُستدعى بـ (من مجلد الدرس):
%    \input{"../../shared/common.tex"}
%
%  يوفّر: التمهيد، الخطوط، الألوان (أزرق/أسود/أبيض)، إطار الصفحة، الترويسة،
%          التذييل (الصفحة س/ص)، شبكة الجداول بخطوط عمودية، صناديق وشارات وأيقونات،
%          وكل ماكروهات الوثائق عبر design.tex + lesson-macros.tex.
% ══════════════════════════════════════════════════════════════════════════════

% ── Geometry ──────────────────────────────────────────────────────────────────
\usepackage[
  top=2.8cm,
  bottom=2.4cm,
  left=2.2cm,
  right=2.2cm,
  headheight=26pt          % MUST be >= 26pt (fancyhdr)
]{geometry}

% ── Packages: fancyhdr and friends BEFORE polyglossia ────────────────────────
\usepackage{fancyhdr}
\usepackage{titlesec}
\usepackage{enumitem}
\usepackage[table]{xcolor}   % table → \rowcolors / \rowcolor / \cellcolor
\usepackage{tikz}
\usetikzlibrary{positioning}
\usetikzlibrary{arrows.meta}
\usetikzlibrary{shapes.geometric}
\usepackage{eso-pic}
\usepackage{tcolorbox}
\tcbuselibrary{skins, breakable}
\usepackage{array}
\usepackage{tabularx}
\usepackage{longtable}
\usepackage{multirow}
\usepackage{lastpage}      % للتذييل: الصفحة س / ص
\usepackage{amsmath}
\usepackage{etoolbox}
\usepackage{hyperref}

% روابط بلا إطار (يمنع المستطيل حول رقم الصفحة في التذييل)
\hypersetup{hidelinks}

% ── Arabic / RTL — load AFTER fancyhdr ───────────────────────────────────────
\usepackage{polyglossia}
\setmainlanguage[numerals=western]{arabic}
\setotherlanguage{english}

% ── Fonts (Amiri — خط النسخ المعتمد) ────────────────────────────────────────
\setmainfont{Amiri}[
  BoldFont = * Bold,
  ItalicFont = * Italic,
]
\newfontfamily\arabicfont{Amiri}[
  Script = Arabic,
  BoldFont = * Bold,
]
\setmonofont{DejaVu Sans Mono}

% ── TYPOGRAPHY: تدرّج طباعي وتهوية ─────────────────────────────────────────
\linespread{1.18}
\setlength{\parindent}{0pt}
\setlength{\parskip}{5pt plus 1pt}
\renewcommand{\arraystretch}{1.2}


% ── مسار مجلد shared (قابل للتجاوز قبل \input{common.tex}) ──────────────────
%   • من lessons/<درس>/  → الافتراضي "../../shared/"
%   • من documents/      → \def\shareddir{../shared/}
\providecommand{\shareddir}{../../shared/}
\makeatletter
\def\input@path{{\shareddir}}
\makeatother

% ── طبقة التصميم (الألوان، الجداول، الصناديق، الشارات، الأيقونات) ────────────
\input{design.tex}

% ── Page Border: 1cm from all four edges (إطار أزرق) ────────────────────────
\AddToShipoutPictureBG{%
  \begin{tikzpicture}[remember picture, overlay]
    \draw[primary, line width=1.1pt]
      ([xshift=-1cm, yshift=-1cm]current page.north east)
      rectangle
      ([xshift=1cm, yshift=1cm]current page.south west);
  \end{tikzpicture}%
}

% ── Header / Footer ──────────────────────────────────────────────────────────
\input{header.tex}
\input{footer.tex}

% ── ماكروهات الوثائق المشتركة ───────────────────────────────────────────────
\input{lesson-macros.tex}

% ── Section formatting ───────────────────────────────────────────────────────
\titleformat{\section}
  {\large\bfseries\color{primary}}
  {\colorbox{primary}{\color{white}\thesection}}
  {0.7em}{}
  [\vspace{-0.15em}{\color{primary}\titlerule[1.2pt]}]

\titleformat{\subsection}
  {\normalsize\bfseries\color{secondary}}{\thesubsection}{0.5em}{}

% عنوان قسم بدون ترقيم — نص بارز فقط (بلا قاعدة)
\newcommand{\plainsec}[1]{%
  \par\vspace{8pt}%
  {\large\bfseries\color{primary} #1}\par\vspace{5pt}%
}

% ── Helpers ──────────────────────────────────────────────────────────────────
% سطر نقاط للإجابة (يجب أن يكون داخل صندوق ليرتسم)
\newcommand{\answerline}{\par\vspace{1pt}\noindent\makebox[\linewidth]{\dotfill}}
% عدّة أسطر إجابة — لا تُستعمل داخل خلية جدول
\newcommand{\answerlines}[1]{%
  \begingroup\count255=0 \loop\ifnum\count255<#1 \answerline\advance\count255 by 1 \repeat\endgroup}
% فراغ للكتابة على السطر
\newcommand{\blank}[1][3cm]{\underline{\hspace{#1}}}
% مربّع اختيار
\newcommand{\cb}{\raisebox{0.15ex}{\fbox{\rule{0pt}{1.1ex}\hspace{1em}}}}
) ──────────────────
%   • من lessons/<درس>/  → الافتراضي "../../shared/"
%   • من documents/      → \def\shareddir{../shared/}
\providecommand{\shareddir}{../../shared/}
\makeatletter
\def\input@path{{\shareddir}}
\makeatother

% ── طبقة التصميم (الألوان، الجداول، الصناديق، الشارات، الأيقونات) ────────────
% ══════════════════════════════════════════════════════════════════════════════
%  design.tex — طبقة التصميم الموحّدة (أزرق / أسود / أبيض فقط)
%  يُستدعى من common.tex (بعد تحميل الحزم) — لا يحمّل أي حزمة هنا.
%
%  يوفّر: الألوان، شبكة الجداول (خطوط عمودية + تظليل متناوب)، العناوين الملوّنة،
%          \sectionbreak (صفحة جديدة قبل كل قسم)، \badge، أيقونات، وصناديق ملاحظات.
% ══════════════════════════════════════════════════════════════════════════════

% ── الألوان (أزرق / أسود / أبيض حصراً) ───────────────────────────────────────
\definecolor{primary}{HTML}{0D47A1}      % أزرق داكن — العناوين ورؤوس الجداول
\definecolor{secondary}{HTML}{1976D2}    % أزرق متوسط — عناوين فرعية
\definecolor{blue}{HTML}{42A5F5}         % أزرق فاتح — شارات ولمسات
\definecolor{darktext}{HTML}{111111}     % أسود للنص
\definecolor{lightbg}{HTML}{E8F0FE}      % خلفية فاتحة جداً
\definecolor{zebra}{HTML}{F2F6FC}        % تظليل الصفوف المتناوب
\definecolor{rulegray}{HTML}{9E9E9E}     % رمادي خطوط الجداول

% تظليل متناوب افتراضي لكل الجداول (أبيض / zebra)
\rowcolors{2}{white}{zebra}

% ── شبكة الجداول: خطوط عمودية + رأس ملوّن + تظليل متناوب ────────────────────
% استعمل قبل كل tabular/x:
%   \gridsetup          — يهيّئ سماكة الخطوط والمسافات
%   \thead{النص}        — خلية رأس بلون primary ونص أبيض
\newcommand{\gridsetup}{%
  \renewcommand{\arraystretch}{1.35}%
  \setlength{\arrayrulewidth}{0.4pt}%
  \arrayrulecolor{rulegray}%
}

% ── فاصل الأقسام: صفحة جديدة + شريط أزرق ممتلئ ──────────────────────────────
% \sectionbreak{العنوان}
\newcommand{\sectionbreak}[1]{%
  \clearpage
  \begin{center}
    \begin{tcolorbox}[enhanced, colback=primary, colframe=primary,
        coltext=white, arc=3pt, boxsep=5pt, left=2pt, right=2pt, top=3pt, bottom=3pt]
      {\LARGE\bfseries #1}
    \end{tcolorbox}
  \end{center}
  \vspace{0.2cm}
}

% ── شارة مرقّمة دائرية ───────────────────────────────────────────────────────
% \badge{الرقم}
\newcommand{\badge}[1]{%
  \tikz[baseline=(b.base)]\node[circle, fill=primary, text=white,
    inner sep=4pt, font=\bfseries] (b) {#1};%
}

% ── أيقونات مرسومة بـ TikZ (لا رموز/إيموجي) ─────────────────────────────────
% علامة صح
\newcommand{\iccheck}{\tikz[baseline=-0.6ex]{%
  \draw[blue, line width=1.4pt, line cap=round] (0,0.12) -- (0.22,0.34) -- (0.6,-0.18);}}
% نجمة
\newcommand{\icstar}{\tikz[baseline=-0.6ex]\node[star, star points=5, star point ratio=2.2,
  fill=blue, draw=blue, minimum size=3.2mm, inner sep=0pt] {};}
% مصباح فكرة
\newcommand{\icbulb}{\tikz[baseline=-0.6ex]{%
  \fill[blue] (0,0.1) circle (2.6mm);
  \draw[blue, line width=0.7pt] (-0.6mm,-0.16) -- (-0.6mm,-0.3) -- (0.6mm,-0.3) -- (0.6mm,-0.16);}}
% شاشة حاسوب
\newcommand{\icmonitor}{\tikz[baseline=-0.6ex]{%
  \draw[blue, line width=0.8pt] (-0.22,-0.05) rectangle (0.22,0.25);
  \draw[blue, line width=0.8pt] (-0.08,-0.05) -- (-0.08,-0.15) -- (0.08,-0.15) -- (0.08,-0.05);
  \draw[blue, line width=0.8pt] (0,-0.15) -- (0,-0.22);}}
% كتاب
\newcommand{\icbook}{\tikz[baseline=-0.6ex]{%
  \draw[blue, line width=0.8pt] (-0.2,-0.2) -- (-0.2,0.18) -- (0,0.22) -- (0.2,0.18) -- (0.2,-0.2);
  \draw[blue, line width=0.8pt] (0,-0.2) -- (0,0.22);}}
% قلم
\newcommand{\icpen}{\tikz[baseline=-0.6ex]{%
  \draw[blue, line width=0.8pt] (-0.24,-0.18) -- (0.16,0.22);
  \draw[blue, fill=white, line width=0.8pt] (-0.24,-0.18) -- (-0.12,-0.24) -- (-0.05,-0.15) -- cycle;}}

% ── صناديق ملاحظات (شريط عنوان أزرق بنص أبيض) ───────────────────────────────
\newtcolorbox{goalbox}[1][]{
  enhanced, breakable, colback=lightbg, colframe=primary,
  colbacktitle=primary, coltitle=white, fonttitle=\bfseries, title=#1,
  boxrule=0pt, arc=3pt, left=11pt, right=11pt, top=6pt, bottom=8pt,
}
\newtcolorbox{activitybox}[1][]{
  enhanced, breakable, colback=white, colframe=secondary,
  colbacktitle=secondary, coltitle=white, fonttitle=\bfseries, title=#1,
  boxrule=0pt, arc=3pt, left=11pt, right=11pt, top=6pt, bottom=8pt,
}
\newtcolorbox{rememberbox}[1][]{
  enhanced, breakable, colback=lightbg, colframe=secondary,
  colbacktitle=secondary, coltitle=white, fonttitle=\bfseries, title=#1,
  boxrule=0pt, arc=3pt, left=11pt, right=11pt, top=6pt, bottom=8pt,
}
\newtcolorbox{tipbox}[1][]{
  enhanced, breakable, colback=white, colframe=blue,
  colbacktitle=blue, coltitle=white, fonttitle=\bfseries, title=#1,
  boxrule=0pt, arc=3pt, left=11pt, right=11pt, top=6pt, bottom=8pt,
}
\newtcolorbox{warnbox}[1][]{
  enhanced, breakable, colback=white, colframe=darktext,
  colbacktitle=darktext, coltitle=white, fonttitle=\bfseries, title=#1,
  boxrule=0pt, arc=3pt, left=11pt, right=11pt, top=6pt, bottom=8pt,
}


% ── Page Border: 1cm from all four edges (إطار أزرق) ────────────────────────
\AddToShipoutPictureBG{%
  \begin{tikzpicture}[remember picture, overlay]
    \draw[primary, line width=1.1pt]
      ([xshift=-1cm, yshift=-1cm]current page.north east)
      rectangle
      ([xshift=1cm, yshift=1cm]current page.south west);
  \end{tikzpicture}%
}

% ── Header / Footer ──────────────────────────────────────────────────────────
% ══════════════════════════════════════════════════════════════════════════════
%  المشترك — الترويسة الرسمية الموحّدة لوثائق الدروس
%  يُستدعى بـ: \input{../_common/header.tex}   (بعد preamble.tex)
%  ثمّ الاستعمال: \officialheader
% ══════════════════════════════════════════════════════════════════════════════

\newcommand{\officialheader}{%
  \par\noindent
  \begin{minipage}{\linewidth}%
    \centering
    {\Large\bfseries\color{primary} الجمهورية الجزائرية الديمقراطية الشعبية}\par
    {\large\bfseries\color{primary} وزارة التربية الوطنية}\par
    \vspace{2pt}%
    \raggedleft
    {\large\color{primary} مديرية التربية لولاية أدرار}\par
    {\large\color{primary} ثانوية الشهيد حكومي العيد أدرار}\par
  \end{minipage}\par
}

\usepackage{lastpage}
\newcommand{\officialfooter}{%
  \pagestyle{fancy}%
  \fancyhf{}%
  \renewcommand{\headrulewidth}{0pt}%
  \renewcommand{\footrulewidth}{0.4pt}%
  \renewcommand{\footrule}{%
    {\color{rulegray}\hrule height \footrulewidth\vskip-\footrulewidth}}%
  \fancyfoot[L]{\footnotesize\color{rulegray} الأستاذ لفقير عبدالجليل}%
  \fancyfoot[R]{\footnotesize\color{rulegray} السنة الدراسية 2027/2026}%
  \fancyfoot[C]{\footnotesize\color{rulegray} الصفحة \thepage\,/\,\pageref{LastPage}}%
}

\newcommand{\pagenum}{%
  \fancyfoot[C]{\footnotesize\color{rulegray} الصفحة \thepage\,/\,\pageref{LastPage}}%
}
\newcommand{\nopagenum}{%
  \fancyfoot[C]{}%
}
\newcommand{\nofooter}{\thispagestyle{empty}}
\newcommand{\nopagefooter}{\thispagestyle{empty}}


% ── ماكروهات الوثائق المشتركة ───────────────────────────────────────────────
% ══════════════════════════════════════════════════════════════════════════════
%  lesson-macros.tex — كل ماكروهات الوثائق مشتركة (شبكة جداول بخطوط عمودية)
%  يُستدعى من common.tex — لا يحمّل أي حزمة هنا.
%
%  ملاحظة محاذاة: لا نستعمل أعمدة p{} بجوار أعمدة X — تختلف نقاط ارتكازهما
%  العمودية (بمقدار \ht\strutbox) فتنزل التسمية عن القيمة. الحل: خلايا التسمية
%  تُغلَّف في \parbox[t] (أو minipage[t]) داخل عمود c، وتتماشى مع عمود X.
% ══════════════════════════════════════════════════════════════════════════════

% ── خلايا محاذاة علوية ───────────────────────────────────────────────────────
% \celllabel{العرض}{المحتوى}   ← تسمية يمين (raggedleft)
% \cellcenter{العرض}{المحتوى}  ← وسط
\newcommand{\celllabel}[2]{\parbox[t]{#1}{\raggedleft #2}}
\newcommand{\cellcenter}[2]{\parbox[t]{#1}{\centering #2}}

% ── خلية رأس: تُستعمل مع \rowcolor{primary} (نص أبيض وسط) ──────────────────
\newcommand{\thead}[2][\linewidth]{\parbox[t]{#1}{\centering\textcolor{white}{\bfseries #2}}}

% ── خلية رأس مبسّطة (خلفية primary ونص أبيض) — للجداول الممتدة/الأفقية ──────
\newcommand{\thd}[1]{\cellcolor{primary}\color{white}\textbf{#1}}

% ── البطاقة الفنية ───────────────────────────────────────────────────────────
\newcommand{\cardsetup}{%
  \renewcommand{\arraystretch}{1.3}%
  \setlength{\arrayrulewidth}{0.4pt}%
  \arrayrulecolor{rulegray}%
}
\newcommand{\cardrow}[2]{\celllabel{3.2cm}{\textbf{\color{primary}#1}} & #2 \\ \hline}
\newcommand{\cardrowtwo}[4]{\celllabel{3.2cm}{\textbf{\color{primary}#1}} & #2 \hfill \textbf{\color{primary}#3:} #4 \\ \hline}
\newcommand{\cardtitle}[1]{%
  \begin{center}
    {\LARGE\bfseries\color{primary} #1}
  \end{center}
  \vspace{6pt}
}

% ── مراحل سير الوضعية التعلمية (مذكرة الدرس) ────────────────────────────────
\newcommand{\stageheader}{%
  \noindent\small
  \begin{tabularx}{\textwidth}{@{}|c|X|c|c|c|c|@{}}
    \hline
    \rowcolor{primary}
    \thead[1.8cm]{مرحلة الدرس} & \thead{مهمات التعلم والهدف منها} &
    \thead[2.7cm]{نشاط الأستاذ} & \thead[2.7cm]{نشاط المتعلم} &
    \thead[2.0cm]{الاستراتيجية} & \thead[1.8cm]{نوع التقويم} \\
    \hline
  \end{tabularx}\par\vspace{-0.4pt}%
}
\newcommand{\stagecontent}[6]{%
  \noindent\small
  \begin{tabularx}{\textwidth}{@{}|c|X|c|c|c|c|@{}}
    \cellcenter{1.8cm}{\textbf{\color{primary}#1}} &
    \celllabel{\linewidth}{#2} &
    \celllabel{2.7cm}{#3} &
    \celllabel{2.7cm}{#4} &
    \cellcenter{2.0cm}{\textbf{\color{secondary}#5}} &
    \cellcenter{1.8cm}{\textbf{#6}} \\
    \hline
  \end{tabularx}\par\vspace{-0.4pt}%
}

% ── مراحل سير الحصة التطبيقية (مذكرة TP) ────────────────────────────────────
\newcommand{\tpstage}[3]{%
  \noindent
  \begin{tabularx}{\textwidth}{@{}|c|c|X|@{}}
    \hline
    \rowcolor{primary}
    \thead[3.0cm]{#1} & \thead[2.0cm]{#2} & \thead{ } \\
    \hline
    \cellcenter{3.0cm}{\textbf{\color{primary}#1}} &
    \cellcenter{2.0cm}{\textbf{#2}} &
    \celllabel{\linewidth}{#3} \\
    \hline
  \end{tabularx}\par\vspace{-0.4pt}%
}

% ── جدول المهام التطبيقية ───────────────────────────────────────────────────
\newcommand{\taskheader}{%
  \noindent
  \begin{tabularx}{\textwidth}{@{}|c|X|X|@{}}
    \hline
    \rowcolor{primary}
    \thead[2.7cm]{المهمة التطبيقية} & \thead{تعليمات الأستاذ والخطوات العملية} &
    \thead{عمل المتعلم على الحاسوب} \\
    \hline
  \end{tabularx}\par\vspace{-0.4pt}%
}
\newcommand{\taskrow}[3]{%
  \noindent
  \begin{tabularx}{\textwidth}{@{}|c|X|X|@{}}
    \cellcenter{2.7cm}{\textbf{\color{primary}#1}} &
    \celllabel{\linewidth}{#2} &
    \celllabel{\linewidth}{#3} \\
    \hline
  \end{tabularx}\par\vspace{-0.4pt}%
}

% ── وثيقة التلميذ ───────────────────────────────────────────────────────────
\newcommand{\sheettitle}[4]{%
  \begin{center}
    {\Large\bfseries\color{primary} وثيقة التلميذ — الأعمال التطبيقية رقم: #1}
  \end{center}
  \vspace{2pt}
  \begin{center}
    {\large\bfseries\color{secondary} #2}
  \end{center}
  \begin{center}
    {\normalsize\color{darktext} المكان: #3 \quad$\cdot$\quad المدة الزمنية: #4}
  \end{center}
  \vspace{6pt}
}
\newcommand{\studentid}{%
  {\centering
  \begin{tabularx}{0.9\textwidth}{@{}|c|X|c|c|@{}}
    \hline
    \celllabel{2.4cm}{\textbf{\color{primary}اللقب والاسم:}} & \blank[4.8cm] &
    \celllabel{2.7cm}{\textbf{\color{primary}القسم / الفوج:}} & \blank[2.2cm] \\
    \hline
    \celllabel{2.4cm}{\textbf{\color{primary}رقم الجهاز:}} & \blank[4.8cm] &
    \celllabel{2.7cm}{\textbf{\color{primary}التاريخ:}} & \blank[2.2cm] \\
    \hline
  \end{tabularx}\par}\vspace{8pt}%
}
\newcommand{\specheader}{%
  \begingroup\centering
  \begin{tabular}{@{}|p{6.0cm}|p{9.0cm}|@{}}
    \hline
    \rowcolor{primary}
    \thead[6.0cm]{العنصر المستهدف} & \thead[9.0cm]{المواصفات المدوّنة على الجهاز} \\
    \hline
}
\newcommand{\specrow}[2]{#1 & \specdots{#2} \\ \hline}
\newcommand{\specend}{\end{tabular}\par\endgroup\vspace{8pt}}
\newcommand{\specdots}[1]{%
  \makebox[\linewidth]{\dotfill}%
  \ifnum#1>1 \\\makebox[\linewidth]{\dotfill}\fi
  \ifnum#1>2 \\\makebox[\linewidth]{\dotfill}\fi
  \ifnum#1>3 \\\makebox[\linewidth]{\dotfill}\fi
}
\newcommand{\graderow}[2]{#1 & #2 \\ \hline}

% ── ملخّص الدرس ─────────────────────────────────────────────────────────────
\newcommand{\summarytitle}[3]{%
  \begin{center}
    {\large\bfseries\color{primary} ملخّص الدرس — الأثر الكتابي على الكراس}
  \end{center}
  \vspace{2pt}
  {\cardsetup
  \centering
  \begin{tabularx}{0.9\textwidth}{@{}|c|X|@{}}
    \hline
    \celllabel{3.0cm}{\textbf{\color{primary}المجال التعلمي}} & #1 \\ \hline
    \celllabel{3.0cm}{\textbf{\color{primary}الوحدة التعلمية}} & #2 \\ \hline
    \celllabel{3.0cm}{\textbf{\color{primary}الموضوع}} & #3 \\ \hline
  \end{tabularx}\par}\vspace{8pt}
}
\newcommand{\summaryaxis}[2]{%
  \par\vspace{6pt}%
  {\normalsize\bfseries\color{primary} #1. #2}\par\vspace{3pt}%
}
\newcommand{\term}[2]{%
  \textbullet\ \textbf{\color{primary}#1:} #2\par\vspace{2pt}%
}

% ── البرنامج السنوي (قائمة، لا جدول) ────────────────────────────────────────
\newcommand{\domain}[2]{%
  \par\vspace{10pt}%
  {\large\bfseries\color{primary} المجال التعلمي #1: #2}\par\vspace{5pt}%
}
\newlist{units}{enumerate}{1}
\setlist[units,1]{%
  label=\textcolor{secondary}{\bfseries\arabic*.},
  labelwidth=1.6em, left=0.4em .. 1.8em,
  itemsep=3pt, topsep=3pt, parsep=0pt,
}


% ── Section formatting ───────────────────────────────────────────────────────
\titleformat{\section}
  {\large\bfseries\color{primary}}
  {\colorbox{primary}{\color{white}\thesection}}
  {0.7em}{}
  [\vspace{-0.15em}{\color{primary}\titlerule[1.2pt]}]

\titleformat{\subsection}
  {\normalsize\bfseries\color{secondary}}{\thesubsection}{0.5em}{}

% عنوان قسم بدون ترقيم — نص بارز فقط (بلا قاعدة)
\newcommand{\plainsec}[1]{%
  \par\vspace{8pt}%
  {\large\bfseries\color{primary} #1}\par\vspace{5pt}%
}

% ── Helpers ──────────────────────────────────────────────────────────────────
% سطر نقاط للإجابة (يجب أن يكون داخل صندوق ليرتسم)
\newcommand{\answerline}{\par\vspace{1pt}\noindent\makebox[\linewidth]{\dotfill}}
% عدّة أسطر إجابة — لا تُستعمل داخل خلية جدول
\newcommand{\answerlines}[1]{%
  \begingroup\count255=0 \loop\ifnum\count255<#1 \answerline\advance\count255 by 1 \repeat\endgroup}
% فراغ للكتابة على السطر
\newcommand{\blank}[1][3cm]{\underline{\hspace{#1}}}
% مربّع اختيار
\newcommand{\cb}{\raisebox{0.15ex}{\fbox{\rule{0pt}{1.1ex}\hspace{1em}}}}
) ──────────────────
%   • من lessons/<درس>/  → الافتراضي "../../shared/"
%   • من documents/      → \def\shareddir{../shared/}
\providecommand{\shareddir}{../../shared/}
\makeatletter
\def\input@path{{\shareddir}}
\makeatother

% ── طبقة التصميم (الألوان، الجداول، الصناديق، الشارات، الأيقونات) ────────────
% ══════════════════════════════════════════════════════════════════════════════
%  design.tex — طبقة التصميم الموحّدة (أزرق / أسود / أبيض فقط)
%  يُستدعى من common.tex (بعد تحميل الحزم) — لا يحمّل أي حزمة هنا.
%
%  يوفّر: الألوان، شبكة الجداول (خطوط عمودية + تظليل متناوب)، العناوين الملوّنة،
%          \sectionbreak (صفحة جديدة قبل كل قسم)، \badge، أيقونات، وصناديق ملاحظات.
% ══════════════════════════════════════════════════════════════════════════════

% ── الألوان (أزرق / أسود / أبيض حصراً) ───────────────────────────────────────
\definecolor{primary}{HTML}{0D47A1}      % أزرق داكن — العناوين ورؤوس الجداول
\definecolor{secondary}{HTML}{1976D2}    % أزرق متوسط — عناوين فرعية
\definecolor{blue}{HTML}{42A5F5}         % أزرق فاتح — شارات ولمسات
\definecolor{darktext}{HTML}{111111}     % أسود للنص
\definecolor{lightbg}{HTML}{E8F0FE}      % خلفية فاتحة جداً
\definecolor{zebra}{HTML}{F2F6FC}        % تظليل الصفوف المتناوب
\definecolor{rulegray}{HTML}{9E9E9E}     % رمادي خطوط الجداول

% تظليل متناوب افتراضي لكل الجداول (أبيض / zebra)
\rowcolors{2}{white}{zebra}

% ── شبكة الجداول: خطوط عمودية + رأس ملوّن + تظليل متناوب ────────────────────
% استعمل قبل كل tabular/x:
%   \gridsetup          — يهيّئ سماكة الخطوط والمسافات
%   \thead{النص}        — خلية رأس بلون primary ونص أبيض
\newcommand{\gridsetup}{%
  \renewcommand{\arraystretch}{1.35}%
  \setlength{\arrayrulewidth}{0.4pt}%
  \arrayrulecolor{rulegray}%
}

% ── فاصل الأقسام: صفحة جديدة + شريط أزرق ممتلئ ──────────────────────────────
% \sectionbreak{العنوان}
\newcommand{\sectionbreak}[1]{%
  \clearpage
  \begin{center}
    \begin{tcolorbox}[enhanced, colback=primary, colframe=primary,
        coltext=white, arc=3pt, boxsep=5pt, left=2pt, right=2pt, top=3pt, bottom=3pt]
      {\LARGE\bfseries #1}
    \end{tcolorbox}
  \end{center}
  \vspace{0.2cm}
}

% ── شارة مرقّمة دائرية ───────────────────────────────────────────────────────
% \badge{الرقم}
\newcommand{\badge}[1]{%
  \tikz[baseline=(b.base)]\node[circle, fill=primary, text=white,
    inner sep=4pt, font=\bfseries] (b) {#1};%
}

% ── أيقونات مرسومة بـ TikZ (لا رموز/إيموجي) ─────────────────────────────────
% علامة صح
\newcommand{\iccheck}{\tikz[baseline=-0.6ex]{%
  \draw[blue, line width=1.4pt, line cap=round] (0,0.12) -- (0.22,0.34) -- (0.6,-0.18);}}
% نجمة
\newcommand{\icstar}{\tikz[baseline=-0.6ex]\node[star, star points=5, star point ratio=2.2,
  fill=blue, draw=blue, minimum size=3.2mm, inner sep=0pt] {};}
% مصباح فكرة
\newcommand{\icbulb}{\tikz[baseline=-0.6ex]{%
  \fill[blue] (0,0.1) circle (2.6mm);
  \draw[blue, line width=0.7pt] (-0.6mm,-0.16) -- (-0.6mm,-0.3) -- (0.6mm,-0.3) -- (0.6mm,-0.16);}}
% شاشة حاسوب
\newcommand{\icmonitor}{\tikz[baseline=-0.6ex]{%
  \draw[blue, line width=0.8pt] (-0.22,-0.05) rectangle (0.22,0.25);
  \draw[blue, line width=0.8pt] (-0.08,-0.05) -- (-0.08,-0.15) -- (0.08,-0.15) -- (0.08,-0.05);
  \draw[blue, line width=0.8pt] (0,-0.15) -- (0,-0.22);}}
% كتاب
\newcommand{\icbook}{\tikz[baseline=-0.6ex]{%
  \draw[blue, line width=0.8pt] (-0.2,-0.2) -- (-0.2,0.18) -- (0,0.22) -- (0.2,0.18) -- (0.2,-0.2);
  \draw[blue, line width=0.8pt] (0,-0.2) -- (0,0.22);}}
% قلم
\newcommand{\icpen}{\tikz[baseline=-0.6ex]{%
  \draw[blue, line width=0.8pt] (-0.24,-0.18) -- (0.16,0.22);
  \draw[blue, fill=white, line width=0.8pt] (-0.24,-0.18) -- (-0.12,-0.24) -- (-0.05,-0.15) -- cycle;}}

% ── صناديق ملاحظات (شريط عنوان أزرق بنص أبيض) ───────────────────────────────
\newtcolorbox{goalbox}[1][]{
  enhanced, breakable, colback=lightbg, colframe=primary,
  colbacktitle=primary, coltitle=white, fonttitle=\bfseries, title=#1,
  boxrule=0pt, arc=3pt, left=11pt, right=11pt, top=6pt, bottom=8pt,
}
\newtcolorbox{activitybox}[1][]{
  enhanced, breakable, colback=white, colframe=secondary,
  colbacktitle=secondary, coltitle=white, fonttitle=\bfseries, title=#1,
  boxrule=0pt, arc=3pt, left=11pt, right=11pt, top=6pt, bottom=8pt,
}
\newtcolorbox{rememberbox}[1][]{
  enhanced, breakable, colback=lightbg, colframe=secondary,
  colbacktitle=secondary, coltitle=white, fonttitle=\bfseries, title=#1,
  boxrule=0pt, arc=3pt, left=11pt, right=11pt, top=6pt, bottom=8pt,
}
\newtcolorbox{tipbox}[1][]{
  enhanced, breakable, colback=white, colframe=blue,
  colbacktitle=blue, coltitle=white, fonttitle=\bfseries, title=#1,
  boxrule=0pt, arc=3pt, left=11pt, right=11pt, top=6pt, bottom=8pt,
}
\newtcolorbox{warnbox}[1][]{
  enhanced, breakable, colback=white, colframe=darktext,
  colbacktitle=darktext, coltitle=white, fonttitle=\bfseries, title=#1,
  boxrule=0pt, arc=3pt, left=11pt, right=11pt, top=6pt, bottom=8pt,
}


% ── Page Border: 1cm from all four edges (إطار أزرق) ────────────────────────
\AddToShipoutPictureBG{%
  \begin{tikzpicture}[remember picture, overlay]
    \draw[primary, line width=1.1pt]
      ([xshift=-1cm, yshift=-1cm]current page.north east)
      rectangle
      ([xshift=1cm, yshift=1cm]current page.south west);
  \end{tikzpicture}%
}

% ── Header / Footer ──────────────────────────────────────────────────────────
% ══════════════════════════════════════════════════════════════════════════════
%  المشترك — الترويسة الرسمية الموحّدة لوثائق الدروس
%  يُستدعى بـ: \input{../_common/header.tex}   (بعد preamble.tex)
%  ثمّ الاستعمال: \officialheader
% ══════════════════════════════════════════════════════════════════════════════

\newcommand{\officialheader}{%
  \par\noindent
  \begin{minipage}{\linewidth}%
    \centering
    {\Large\bfseries\color{primary} الجمهورية الجزائرية الديمقراطية الشعبية}\par
    {\large\bfseries\color{primary} وزارة التربية الوطنية}\par
    \vspace{2pt}%
    \raggedleft
    {\large\color{primary} مديرية التربية لولاية أدرار}\par
    {\large\color{primary} ثانوية الشهيد حكومي العيد أدرار}\par
  \end{minipage}\par
}

\usepackage{lastpage}
\newcommand{\officialfooter}{%
  \pagestyle{fancy}%
  \fancyhf{}%
  \renewcommand{\headrulewidth}{0pt}%
  \renewcommand{\footrulewidth}{0.4pt}%
  \renewcommand{\footrule}{%
    {\color{rulegray}\hrule height \footrulewidth\vskip-\footrulewidth}}%
  \fancyfoot[L]{\footnotesize\color{rulegray} الأستاذ لفقير عبدالجليل}%
  \fancyfoot[R]{\footnotesize\color{rulegray} السنة الدراسية 2027/2026}%
  \fancyfoot[C]{\footnotesize\color{rulegray} الصفحة \thepage\,/\,\pageref{LastPage}}%
}

\newcommand{\pagenum}{%
  \fancyfoot[C]{\footnotesize\color{rulegray} الصفحة \thepage\,/\,\pageref{LastPage}}%
}
\newcommand{\nopagenum}{%
  \fancyfoot[C]{}%
}
\newcommand{\nofooter}{\thispagestyle{empty}}
\newcommand{\nopagefooter}{\thispagestyle{empty}}


% ── ماكروهات الوثائق المشتركة ───────────────────────────────────────────────
% ══════════════════════════════════════════════════════════════════════════════
%  lesson-macros.tex — كل ماكروهات الوثائق مشتركة (شبكة جداول بخطوط عمودية)
%  يُستدعى من common.tex — لا يحمّل أي حزمة هنا.
%
%  ملاحظة محاذاة: لا نستعمل أعمدة p{} بجوار أعمدة X — تختلف نقاط ارتكازهما
%  العمودية (بمقدار \ht\strutbox) فتنزل التسمية عن القيمة. الحل: خلايا التسمية
%  تُغلَّف في \parbox[t] (أو minipage[t]) داخل عمود c، وتتماشى مع عمود X.
% ══════════════════════════════════════════════════════════════════════════════

% ── خلايا محاذاة علوية ───────────────────────────────────────────────────────
% \celllabel{العرض}{المحتوى}   ← تسمية يمين (raggedleft)
% \cellcenter{العرض}{المحتوى}  ← وسط
\newcommand{\celllabel}[2]{\parbox[t]{#1}{\raggedleft #2}}
\newcommand{\cellcenter}[2]{\parbox[t]{#1}{\centering #2}}

% ── خلية رأس: تُستعمل مع \rowcolor{primary} (نص أبيض وسط) ──────────────────
\newcommand{\thead}[2][\linewidth]{\parbox[t]{#1}{\centering\textcolor{white}{\bfseries #2}}}

% ── خلية رأس مبسّطة (خلفية primary ونص أبيض) — للجداول الممتدة/الأفقية ──────
\newcommand{\thd}[1]{\cellcolor{primary}\color{white}\textbf{#1}}

% ── البطاقة الفنية ───────────────────────────────────────────────────────────
\newcommand{\cardsetup}{%
  \renewcommand{\arraystretch}{1.3}%
  \setlength{\arrayrulewidth}{0.4pt}%
  \arrayrulecolor{rulegray}%
}
\newcommand{\cardrow}[2]{\celllabel{3.2cm}{\textbf{\color{primary}#1}} & #2 \\ \hline}
\newcommand{\cardrowtwo}[4]{\celllabel{3.2cm}{\textbf{\color{primary}#1}} & #2 \hfill \textbf{\color{primary}#3:} #4 \\ \hline}
\newcommand{\cardtitle}[1]{%
  \begin{center}
    {\LARGE\bfseries\color{primary} #1}
  \end{center}
  \vspace{6pt}
}

% ── مراحل سير الوضعية التعلمية (مذكرة الدرس) ────────────────────────────────
\newcommand{\stageheader}{%
  \noindent\small
  \begin{tabularx}{\textwidth}{@{}|c|X|c|c|c|c|@{}}
    \hline
    \rowcolor{primary}
    \thead[1.8cm]{مرحلة الدرس} & \thead{مهمات التعلم والهدف منها} &
    \thead[2.7cm]{نشاط الأستاذ} & \thead[2.7cm]{نشاط المتعلم} &
    \thead[2.0cm]{الاستراتيجية} & \thead[1.8cm]{نوع التقويم} \\
    \hline
  \end{tabularx}\par\vspace{-0.4pt}%
}
\newcommand{\stagecontent}[6]{%
  \noindent\small
  \begin{tabularx}{\textwidth}{@{}|c|X|c|c|c|c|@{}}
    \cellcenter{1.8cm}{\textbf{\color{primary}#1}} &
    \celllabel{\linewidth}{#2} &
    \celllabel{2.7cm}{#3} &
    \celllabel{2.7cm}{#4} &
    \cellcenter{2.0cm}{\textbf{\color{secondary}#5}} &
    \cellcenter{1.8cm}{\textbf{#6}} \\
    \hline
  \end{tabularx}\par\vspace{-0.4pt}%
}

% ── مراحل سير الحصة التطبيقية (مذكرة TP) ────────────────────────────────────
\newcommand{\tpstage}[3]{%
  \noindent
  \begin{tabularx}{\textwidth}{@{}|c|c|X|@{}}
    \hline
    \rowcolor{primary}
    \thead[3.0cm]{#1} & \thead[2.0cm]{#2} & \thead{ } \\
    \hline
    \cellcenter{3.0cm}{\textbf{\color{primary}#1}} &
    \cellcenter{2.0cm}{\textbf{#2}} &
    \celllabel{\linewidth}{#3} \\
    \hline
  \end{tabularx}\par\vspace{-0.4pt}%
}

% ── جدول المهام التطبيقية ───────────────────────────────────────────────────
\newcommand{\taskheader}{%
  \noindent
  \begin{tabularx}{\textwidth}{@{}|c|X|X|@{}}
    \hline
    \rowcolor{primary}
    \thead[2.7cm]{المهمة التطبيقية} & \thead{تعليمات الأستاذ والخطوات العملية} &
    \thead{عمل المتعلم على الحاسوب} \\
    \hline
  \end{tabularx}\par\vspace{-0.4pt}%
}
\newcommand{\taskrow}[3]{%
  \noindent
  \begin{tabularx}{\textwidth}{@{}|c|X|X|@{}}
    \cellcenter{2.7cm}{\textbf{\color{primary}#1}} &
    \celllabel{\linewidth}{#2} &
    \celllabel{\linewidth}{#3} \\
    \hline
  \end{tabularx}\par\vspace{-0.4pt}%
}

% ── وثيقة التلميذ ───────────────────────────────────────────────────────────
\newcommand{\sheettitle}[4]{%
  \begin{center}
    {\Large\bfseries\color{primary} وثيقة التلميذ — الأعمال التطبيقية رقم: #1}
  \end{center}
  \vspace{2pt}
  \begin{center}
    {\large\bfseries\color{secondary} #2}
  \end{center}
  \begin{center}
    {\normalsize\color{darktext} المكان: #3 \quad$\cdot$\quad المدة الزمنية: #4}
  \end{center}
  \vspace{6pt}
}
\newcommand{\studentid}{%
  {\centering
  \begin{tabularx}{0.9\textwidth}{@{}|c|X|c|c|@{}}
    \hline
    \celllabel{2.4cm}{\textbf{\color{primary}اللقب والاسم:}} & \blank[4.8cm] &
    \celllabel{2.7cm}{\textbf{\color{primary}القسم / الفوج:}} & \blank[2.2cm] \\
    \hline
    \celllabel{2.4cm}{\textbf{\color{primary}رقم الجهاز:}} & \blank[4.8cm] &
    \celllabel{2.7cm}{\textbf{\color{primary}التاريخ:}} & \blank[2.2cm] \\
    \hline
  \end{tabularx}\par}\vspace{8pt}%
}
\newcommand{\specheader}{%
  \begingroup\centering
  \begin{tabular}{@{}|p{6.0cm}|p{9.0cm}|@{}}
    \hline
    \rowcolor{primary}
    \thead[6.0cm]{العنصر المستهدف} & \thead[9.0cm]{المواصفات المدوّنة على الجهاز} \\
    \hline
}
\newcommand{\specrow}[2]{#1 & \specdots{#2} \\ \hline}
\newcommand{\specend}{\end{tabular}\par\endgroup\vspace{8pt}}
\newcommand{\specdots}[1]{%
  \makebox[\linewidth]{\dotfill}%
  \ifnum#1>1 \\\makebox[\linewidth]{\dotfill}\fi
  \ifnum#1>2 \\\makebox[\linewidth]{\dotfill}\fi
  \ifnum#1>3 \\\makebox[\linewidth]{\dotfill}\fi
}
\newcommand{\graderow}[2]{#1 & #2 \\ \hline}

% ── ملخّص الدرس ─────────────────────────────────────────────────────────────
\newcommand{\summarytitle}[3]{%
  \begin{center}
    {\large\bfseries\color{primary} ملخّص الدرس — الأثر الكتابي على الكراس}
  \end{center}
  \vspace{2pt}
  {\cardsetup
  \centering
  \begin{tabularx}{0.9\textwidth}{@{}|c|X|@{}}
    \hline
    \celllabel{3.0cm}{\textbf{\color{primary}المجال التعلمي}} & #1 \\ \hline
    \celllabel{3.0cm}{\textbf{\color{primary}الوحدة التعلمية}} & #2 \\ \hline
    \celllabel{3.0cm}{\textbf{\color{primary}الموضوع}} & #3 \\ \hline
  \end{tabularx}\par}\vspace{8pt}
}
\newcommand{\summaryaxis}[2]{%
  \par\vspace{6pt}%
  {\normalsize\bfseries\color{primary} #1. #2}\par\vspace{3pt}%
}
\newcommand{\term}[2]{%
  \textbullet\ \textbf{\color{primary}#1:} #2\par\vspace{2pt}%
}

% ── البرنامج السنوي (قائمة، لا جدول) ────────────────────────────────────────
\newcommand{\domain}[2]{%
  \par\vspace{10pt}%
  {\large\bfseries\color{primary} المجال التعلمي #1: #2}\par\vspace{5pt}%
}
\newlist{units}{enumerate}{1}
\setlist[units,1]{%
  label=\textcolor{secondary}{\bfseries\arabic*.},
  labelwidth=1.6em, left=0.4em .. 1.8em,
  itemsep=3pt, topsep=3pt, parsep=0pt,
}


% ── Section formatting ───────────────────────────────────────────────────────
\titleformat{\section}
  {\large\bfseries\color{primary}}
  {\colorbox{primary}{\color{white}\thesection}}
  {0.7em}{}
  [\vspace{-0.15em}{\color{primary}\titlerule[1.2pt]}]

\titleformat{\subsection}
  {\normalsize\bfseries\color{secondary}}{\thesubsection}{0.5em}{}

% عنوان قسم بدون ترقيم — نص بارز فقط (بلا قاعدة)
\newcommand{\plainsec}[1]{%
  \par\vspace{8pt}%
  {\large\bfseries\color{primary} #1}\par\vspace{5pt}%
}

% ── Helpers ──────────────────────────────────────────────────────────────────
% سطر نقاط للإجابة (يجب أن يكون داخل صندوق ليرتسم)
\newcommand{\answerline}{\par\vspace{1pt}\noindent\makebox[\linewidth]{\dotfill}}
% عدّة أسطر إجابة — لا تُستعمل داخل خلية جدول
\newcommand{\answerlines}[1]{%
  \begingroup\count255=0 \loop\ifnum\count255<#1 \answerline\advance\count255 by 1 \repeat\endgroup}
% فراغ للكتابة على السطر
\newcommand{\blank}[1][3cm]{\underline{\hspace{#1}}}
% مربّع اختيار
\newcommand{\cb}{\raisebox{0.15ex}{\fbox{\rule{0pt}{1.1ex}\hspace{1em}}}}
) ──────────────────
%   • من lessons/<درس>/  → الافتراضي "../../shared/"
%   • من documents/      → \def\shareddir{../shared/}
\providecommand{\shareddir}{../../shared/}
\makeatletter
\def\input@path{{\shareddir}}
\makeatother

% ── طبقة التصميم (الألوان، الجداول، الصناديق، الشارات، الأيقونات) ────────────
% ══════════════════════════════════════════════════════════════════════════════
%  design.tex — طبقة التصميم الموحّدة (أزرق / أسود / أبيض فقط)
%  يُستدعى من common.tex (بعد تحميل الحزم) — لا يحمّل أي حزمة هنا.
%
%  يوفّر: الألوان، شبكة الجداول (خطوط عمودية + تظليل متناوب)، العناوين الملوّنة،
%          \sectionbreak (صفحة جديدة قبل كل قسم)، \badge، أيقونات، وصناديق ملاحظات.
% ══════════════════════════════════════════════════════════════════════════════

% ── الألوان (أزرق / أسود / أبيض حصراً) ───────────────────────────────────────
\definecolor{primary}{HTML}{0D47A1}      % أزرق داكن — العناوين ورؤوس الجداول
\definecolor{secondary}{HTML}{1976D2}    % أزرق متوسط — عناوين فرعية
\definecolor{blue}{HTML}{42A5F5}         % أزرق فاتح — شارات ولمسات
\definecolor{darktext}{HTML}{111111}     % أسود للنص
\definecolor{lightbg}{HTML}{E8F0FE}      % خلفية فاتحة جداً
\definecolor{zebra}{HTML}{F2F6FC}        % تظليل الصفوف المتناوب
\definecolor{rulegray}{HTML}{9E9E9E}     % رمادي خطوط الجداول

% تظليل متناوب افتراضي لكل الجداول (أبيض / zebra)
\rowcolors{2}{white}{zebra}

% ── شبكة الجداول: خطوط عمودية + رأس ملوّن + تظليل متناوب ────────────────────
% استعمل قبل كل tabular/x:
%   \gridsetup          — يهيّئ سماكة الخطوط والمسافات
%   \thead{النص}        — خلية رأس بلون primary ونص أبيض
\newcommand{\gridsetup}{%
  \renewcommand{\arraystretch}{1.35}%
  \setlength{\arrayrulewidth}{0.4pt}%
  \arrayrulecolor{rulegray}%
}

% ── فاصل الأقسام: صفحة جديدة + شريط أزرق ممتلئ ──────────────────────────────
% \sectionbreak{العنوان}
\newcommand{\sectionbreak}[1]{%
  \clearpage
  \begin{center}
    \begin{tcolorbox}[enhanced, colback=primary, colframe=primary,
        coltext=white, arc=3pt, boxsep=5pt, left=2pt, right=2pt, top=3pt, bottom=3pt]
      {\LARGE\bfseries #1}
    \end{tcolorbox}
  \end{center}
  \vspace{0.2cm}
}

% ── شارة مرقّمة دائرية ───────────────────────────────────────────────────────
% \badge{الرقم}
\newcommand{\badge}[1]{%
  \tikz[baseline=(b.base)]\node[circle, fill=primary, text=white,
    inner sep=4pt, font=\bfseries] (b) {#1};%
}

% ── أيقونات مرسومة بـ TikZ (لا رموز/إيموجي) ─────────────────────────────────
% علامة صح
\newcommand{\iccheck}{\tikz[baseline=-0.6ex]{%
  \draw[blue, line width=1.4pt, line cap=round] (0,0.12) -- (0.22,0.34) -- (0.6,-0.18);}}
% نجمة
\newcommand{\icstar}{\tikz[baseline=-0.6ex]\node[star, star points=5, star point ratio=2.2,
  fill=blue, draw=blue, minimum size=3.2mm, inner sep=0pt] {};}
% مصباح فكرة
\newcommand{\icbulb}{\tikz[baseline=-0.6ex]{%
  \fill[blue] (0,0.1) circle (2.6mm);
  \draw[blue, line width=0.7pt] (-0.6mm,-0.16) -- (-0.6mm,-0.3) -- (0.6mm,-0.3) -- (0.6mm,-0.16);}}
% شاشة حاسوب
\newcommand{\icmonitor}{\tikz[baseline=-0.6ex]{%
  \draw[blue, line width=0.8pt] (-0.22,-0.05) rectangle (0.22,0.25);
  \draw[blue, line width=0.8pt] (-0.08,-0.05) -- (-0.08,-0.15) -- (0.08,-0.15) -- (0.08,-0.05);
  \draw[blue, line width=0.8pt] (0,-0.15) -- (0,-0.22);}}
% كتاب
\newcommand{\icbook}{\tikz[baseline=-0.6ex]{%
  \draw[blue, line width=0.8pt] (-0.2,-0.2) -- (-0.2,0.18) -- (0,0.22) -- (0.2,0.18) -- (0.2,-0.2);
  \draw[blue, line width=0.8pt] (0,-0.2) -- (0,0.22);}}
% قلم
\newcommand{\icpen}{\tikz[baseline=-0.6ex]{%
  \draw[blue, line width=0.8pt] (-0.24,-0.18) -- (0.16,0.22);
  \draw[blue, fill=white, line width=0.8pt] (-0.24,-0.18) -- (-0.12,-0.24) -- (-0.05,-0.15) -- cycle;}}

% ── صناديق ملاحظات (شريط عنوان أزرق بنص أبيض) ───────────────────────────────
\newtcolorbox{goalbox}[1][]{
  enhanced, breakable, colback=lightbg, colframe=primary,
  colbacktitle=primary, coltitle=white, fonttitle=\bfseries, title=#1,
  boxrule=0pt, arc=3pt, left=11pt, right=11pt, top=6pt, bottom=8pt,
}
\newtcolorbox{activitybox}[1][]{
  enhanced, breakable, colback=white, colframe=secondary,
  colbacktitle=secondary, coltitle=white, fonttitle=\bfseries, title=#1,
  boxrule=0pt, arc=3pt, left=11pt, right=11pt, top=6pt, bottom=8pt,
}
\newtcolorbox{rememberbox}[1][]{
  enhanced, breakable, colback=lightbg, colframe=secondary,
  colbacktitle=secondary, coltitle=white, fonttitle=\bfseries, title=#1,
  boxrule=0pt, arc=3pt, left=11pt, right=11pt, top=6pt, bottom=8pt,
}
\newtcolorbox{tipbox}[1][]{
  enhanced, breakable, colback=white, colframe=blue,
  colbacktitle=blue, coltitle=white, fonttitle=\bfseries, title=#1,
  boxrule=0pt, arc=3pt, left=11pt, right=11pt, top=6pt, bottom=8pt,
}
\newtcolorbox{warnbox}[1][]{
  enhanced, breakable, colback=white, colframe=darktext,
  colbacktitle=darktext, coltitle=white, fonttitle=\bfseries, title=#1,
  boxrule=0pt, arc=3pt, left=11pt, right=11pt, top=6pt, bottom=8pt,
}


% ── Page Border: 1cm from all four edges (إطار أزرق) ────────────────────────
\AddToShipoutPictureBG{%
  \begin{tikzpicture}[remember picture, overlay]
    \draw[primary, line width=1.1pt]
      ([xshift=-1cm, yshift=-1cm]current page.north east)
      rectangle
      ([xshift=1cm, yshift=1cm]current page.south west);
  \end{tikzpicture}%
}

% ── Header / Footer ──────────────────────────────────────────────────────────
% ══════════════════════════════════════════════════════════════════════════════
%  المشترك — الترويسة الرسمية الموحّدة لوثائق الدروس
%  يُستدعى بـ: \input{../_common/header.tex}   (بعد preamble.tex)
%  ثمّ الاستعمال: \officialheader
% ══════════════════════════════════════════════════════════════════════════════

\newcommand{\officialheader}{%
  \par\noindent
  \begin{minipage}{\linewidth}%
    \centering
    {\Large\bfseries\color{primary} الجمهورية الجزائرية الديمقراطية الشعبية}\par
    {\large\bfseries\color{primary} وزارة التربية الوطنية}\par
    \vspace{2pt}%
    \raggedleft
    {\large\color{primary} مديرية التربية لولاية أدرار}\par
    {\large\color{primary} ثانوية الشهيد حكومي العيد أدرار}\par
  \end{minipage}\par
}

\usepackage{lastpage}
\newcommand{\officialfooter}{%
  \pagestyle{fancy}%
  \fancyhf{}%
  \renewcommand{\headrulewidth}{0pt}%
  \renewcommand{\footrulewidth}{0.4pt}%
  \renewcommand{\footrule}{%
    {\color{rulegray}\hrule height \footrulewidth\vskip-\footrulewidth}}%
  \fancyfoot[L]{\footnotesize\color{rulegray} الأستاذ لفقير عبدالجليل}%
  \fancyfoot[R]{\footnotesize\color{rulegray} السنة الدراسية 2027/2026}%
  \fancyfoot[C]{\footnotesize\color{rulegray} الصفحة \thepage\,/\,\pageref{LastPage}}%
}

\newcommand{\pagenum}{%
  \fancyfoot[C]{\footnotesize\color{rulegray} الصفحة \thepage\,/\,\pageref{LastPage}}%
}
\newcommand{\nopagenum}{%
  \fancyfoot[C]{}%
}
\newcommand{\nofooter}{\thispagestyle{empty}}
\newcommand{\nopagefooter}{\thispagestyle{empty}}


% ── ماكروهات الوثائق المشتركة ───────────────────────────────────────────────
% ══════════════════════════════════════════════════════════════════════════════
%  lesson-macros.tex — كل ماكروهات الوثائق مشتركة (شبكة جداول بخطوط عمودية)
%  يُستدعى من common.tex — لا يحمّل أي حزمة هنا.
%
%  ملاحظة محاذاة: لا نستعمل أعمدة p{} بجوار أعمدة X — تختلف نقاط ارتكازهما
%  العمودية (بمقدار \ht\strutbox) فتنزل التسمية عن القيمة. الحل: خلايا التسمية
%  تُغلَّف في \parbox[t] (أو minipage[t]) داخل عمود c، وتتماشى مع عمود X.
% ══════════════════════════════════════════════════════════════════════════════

% ── خلايا محاذاة علوية ───────────────────────────────────────────────────────
% \celllabel{العرض}{المحتوى}   ← تسمية يمين (raggedleft)
% \cellcenter{العرض}{المحتوى}  ← وسط
\newcommand{\celllabel}[2]{\parbox[t]{#1}{\raggedleft #2}}
\newcommand{\cellcenter}[2]{\parbox[t]{#1}{\centering #2}}

% ── خلية رأس: تُستعمل مع \rowcolor{primary} (نص أبيض وسط) ──────────────────
\newcommand{\thead}[2][\linewidth]{\parbox[t]{#1}{\centering\textcolor{white}{\bfseries #2}}}

% ── خلية رأس مبسّطة (خلفية primary ونص أبيض) — للجداول الممتدة/الأفقية ──────
\newcommand{\thd}[1]{\cellcolor{primary}\color{white}\textbf{#1}}

% ── البطاقة الفنية ───────────────────────────────────────────────────────────
\newcommand{\cardsetup}{%
  \renewcommand{\arraystretch}{1.3}%
  \setlength{\arrayrulewidth}{0.4pt}%
  \arrayrulecolor{rulegray}%
}
\newcommand{\cardrow}[2]{\celllabel{3.2cm}{\textbf{\color{primary}#1}} & #2 \\ \hline}
\newcommand{\cardrowtwo}[4]{\celllabel{3.2cm}{\textbf{\color{primary}#1}} & #2 \hfill \textbf{\color{primary}#3:} #4 \\ \hline}
\newcommand{\cardtitle}[1]{%
  \begin{center}
    {\LARGE\bfseries\color{primary} #1}
  \end{center}
  \vspace{6pt}
}

% ── مراحل سير الوضعية التعلمية (مذكرة الدرس) ────────────────────────────────
\newcommand{\stageheader}{%
  \noindent\small
  \begin{tabularx}{\textwidth}{@{}|c|X|c|c|c|c|@{}}
    \hline
    \rowcolor{primary}
    \thead[1.8cm]{مرحلة الدرس} & \thead{مهمات التعلم والهدف منها} &
    \thead[2.7cm]{نشاط الأستاذ} & \thead[2.7cm]{نشاط المتعلم} &
    \thead[2.0cm]{الاستراتيجية} & \thead[1.8cm]{نوع التقويم} \\
    \hline
  \end{tabularx}\par\vspace{-0.4pt}%
}
\newcommand{\stagecontent}[6]{%
  \noindent\small
  \begin{tabularx}{\textwidth}{@{}|c|X|c|c|c|c|@{}}
    \cellcenter{1.8cm}{\textbf{\color{primary}#1}} &
    \celllabel{\linewidth}{#2} &
    \celllabel{2.7cm}{#3} &
    \celllabel{2.7cm}{#4} &
    \cellcenter{2.0cm}{\textbf{\color{secondary}#5}} &
    \cellcenter{1.8cm}{\textbf{#6}} \\
    \hline
  \end{tabularx}\par\vspace{-0.4pt}%
}

% ── مراحل سير الحصة التطبيقية (مذكرة TP) ────────────────────────────────────
\newcommand{\tpstage}[3]{%
  \noindent
  \begin{tabularx}{\textwidth}{@{}|c|c|X|@{}}
    \hline
    \rowcolor{primary}
    \thead[3.0cm]{#1} & \thead[2.0cm]{#2} & \thead{ } \\
    \hline
    \cellcenter{3.0cm}{\textbf{\color{primary}#1}} &
    \cellcenter{2.0cm}{\textbf{#2}} &
    \celllabel{\linewidth}{#3} \\
    \hline
  \end{tabularx}\par\vspace{-0.4pt}%
}

% ── جدول المهام التطبيقية ───────────────────────────────────────────────────
\newcommand{\taskheader}{%
  \noindent
  \begin{tabularx}{\textwidth}{@{}|c|X|X|@{}}
    \hline
    \rowcolor{primary}
    \thead[2.7cm]{المهمة التطبيقية} & \thead{تعليمات الأستاذ والخطوات العملية} &
    \thead{عمل المتعلم على الحاسوب} \\
    \hline
  \end{tabularx}\par\vspace{-0.4pt}%
}
\newcommand{\taskrow}[3]{%
  \noindent
  \begin{tabularx}{\textwidth}{@{}|c|X|X|@{}}
    \cellcenter{2.7cm}{\textbf{\color{primary}#1}} &
    \celllabel{\linewidth}{#2} &
    \celllabel{\linewidth}{#3} \\
    \hline
  \end{tabularx}\par\vspace{-0.4pt}%
}

% ── وثيقة التلميذ ───────────────────────────────────────────────────────────
\newcommand{\sheettitle}[4]{%
  \begin{center}
    {\Large\bfseries\color{primary} وثيقة التلميذ — الأعمال التطبيقية رقم: #1}
  \end{center}
  \vspace{2pt}
  \begin{center}
    {\large\bfseries\color{secondary} #2}
  \end{center}
  \begin{center}
    {\normalsize\color{darktext} المكان: #3 \quad$\cdot$\quad المدة الزمنية: #4}
  \end{center}
  \vspace{6pt}
}
\newcommand{\studentid}{%
  {\centering
  \begin{tabularx}{0.9\textwidth}{@{}|c|X|c|c|@{}}
    \hline
    \celllabel{2.4cm}{\textbf{\color{primary}اللقب والاسم:}} & \blank[4.8cm] &
    \celllabel{2.7cm}{\textbf{\color{primary}القسم / الفوج:}} & \blank[2.2cm] \\
    \hline
    \celllabel{2.4cm}{\textbf{\color{primary}رقم الجهاز:}} & \blank[4.8cm] &
    \celllabel{2.7cm}{\textbf{\color{primary}التاريخ:}} & \blank[2.2cm] \\
    \hline
  \end{tabularx}\par}\vspace{8pt}%
}
\newcommand{\specheader}{%
  \begingroup\centering
  \begin{tabular}{@{}|p{6.0cm}|p{9.0cm}|@{}}
    \hline
    \rowcolor{primary}
    \thead[6.0cm]{العنصر المستهدف} & \thead[9.0cm]{المواصفات المدوّنة على الجهاز} \\
    \hline
}
\newcommand{\specrow}[2]{#1 & \specdots{#2} \\ \hline}
\newcommand{\specend}{\end{tabular}\par\endgroup\vspace{8pt}}
\newcommand{\specdots}[1]{%
  \makebox[\linewidth]{\dotfill}%
  \ifnum#1>1 \\\makebox[\linewidth]{\dotfill}\fi
  \ifnum#1>2 \\\makebox[\linewidth]{\dotfill}\fi
  \ifnum#1>3 \\\makebox[\linewidth]{\dotfill}\fi
}
\newcommand{\graderow}[2]{#1 & #2 \\ \hline}

% ── ملخّص الدرس ─────────────────────────────────────────────────────────────
\newcommand{\summarytitle}[3]{%
  \begin{center}
    {\large\bfseries\color{primary} ملخّص الدرس — الأثر الكتابي على الكراس}
  \end{center}
  \vspace{2pt}
  {\cardsetup
  \centering
  \begin{tabularx}{0.9\textwidth}{@{}|c|X|@{}}
    \hline
    \celllabel{3.0cm}{\textbf{\color{primary}المجال التعلمي}} & #1 \\ \hline
    \celllabel{3.0cm}{\textbf{\color{primary}الوحدة التعلمية}} & #2 \\ \hline
    \celllabel{3.0cm}{\textbf{\color{primary}الموضوع}} & #3 \\ \hline
  \end{tabularx}\par}\vspace{8pt}
}
\newcommand{\summaryaxis}[2]{%
  \par\vspace{6pt}%
  {\normalsize\bfseries\color{primary} #1. #2}\par\vspace{3pt}%
}
\newcommand{\term}[2]{%
  \textbullet\ \textbf{\color{primary}#1:} #2\par\vspace{2pt}%
}

% ── البرنامج السنوي (قائمة، لا جدول) ────────────────────────────────────────
\newcommand{\domain}[2]{%
  \par\vspace{10pt}%
  {\large\bfseries\color{primary} المجال التعلمي #1: #2}\par\vspace{5pt}%
}
\newlist{units}{enumerate}{1}
\setlist[units,1]{%
  label=\textcolor{secondary}{\bfseries\arabic*.},
  labelwidth=1.6em, left=0.4em .. 1.8em,
  itemsep=3pt, topsep=3pt, parsep=0pt,
}


% ── Section formatting ───────────────────────────────────────────────────────
\titleformat{\section}
  {\large\bfseries\color{primary}}
  {\colorbox{primary}{\color{white}\thesection}}
  {0.7em}{}
  [\vspace{-0.15em}{\color{primary}\titlerule[1.2pt]}]

\titleformat{\subsection}
  {\normalsize\bfseries\color{secondary}}{\thesubsection}{0.5em}{}

% عنوان قسم بدون ترقيم — نص بارز فقط (بلا قاعدة)
\newcommand{\plainsec}[1]{%
  \par\vspace{8pt}%
  {\large\bfseries\color{primary} #1}\par\vspace{5pt}%
}

% ── Helpers ──────────────────────────────────────────────────────────────────
% سطر نقاط للإجابة (يجب أن يكون داخل صندوق ليرتسم)
\newcommand{\answerline}{\par\vspace{1pt}\noindent\makebox[\linewidth]{\dotfill}}
% عدّة أسطر إجابة — لا تُستعمل داخل خلية جدول
\newcommand{\answerlines}[1]{%
  \begingroup\count255=0 \loop\ifnum\count255<#1 \answerline\advance\count255 by 1 \repeat\endgroup}
% فراغ للكتابة على السطر
\newcommand{\blank}[1][3cm]{\underline{\hspace{#1}}}
% مربّع اختيار
\newcommand{\cb}{\raisebox{0.15ex}{\fbox{\rule{0pt}{1.1ex}\hspace{1em}}}}
